\documentclass[11pt,a4paper]{article}
\usepackage[top=1.7cm,bottom=1.5cm,left=1.5cm,right=1.5cm,headsep=8pt,footskip=16pt]{geometry}
\usepackage{enumitem}
\usepackage{amssymb}
\usepackage{array}
\usepackage[table]{xcolor}
\usepackage{fancyhdr}
\usepackage{lastpage}
\usepackage{multicol}
\usepackage{polyglossia}
\setdefaultlanguage[calendar=gregorian, numerals=maghrib]{arabic}
\setotherlanguage{english}
\IfFontExistsTF{Amiri}{%
  \newfontfamily\arabicfont[Script=Arabic, Scale=1.08]{Amiri}
}{%
  \newfontfamily\arabicfont[Script=Arabic, Scale=1.08, Path=../fonts/, Extension=.ttf,
    UprightFont=*-Regular, BoldFont=*-Bold, ItalicFont=*-Italic,
    BoldItalicFont=*-BoldItalic]{Amiri}
}
\IfFontExistsTF{AlBattar}{%
  \newfontfamily\khatt[Script=Arabic, Scale=1.15]{AlBattar}
}{%
  \newfontfamily\khatt[Script=Arabic, Scale=1.15, Path=../fonts/,
    BoldFont=ae_AlBattar.ttf]{ae_AlBattar.ttf}
}
\renewcommand{\headrulewidth}{0.8pt}
\renewcommand{\footrulewidth}{0.8pt}
\fancyhf{}
\rhead{\khatt المدرسة العليا للأساتذة ببوزريعة}
\lhead{\khatt ملحقة بومرداس}
\lfoot{\khatt الأستاذ: د. بن عبد الرزاق يوسف}
\cfoot{\thepage/\pageref{LastPage}}
\pagestyle{fancy}
\parindent 0pt
\setlength{\columnsep}{0.6cm}
\setlength{\columnseprule}{0.4pt}
\setlist[enumerate,1]{label=\textbf{\arabic*.}, leftmargin=*, itemsep=5pt, topsep=2pt}
\begin{document}

\linespread{1.08}\selectfont
\begin{center}
{\LARGE\khatt\textbf{الاختبار التاسع في الذكاء الاصطناعي}}\\[4pt]
{\normalsize وضعيات إشكالية ــ أسئلة متعددة الاختيارات ــ الفصول من 1 إلى 6}
\end{center}
\vspace{-2pt}\rule{\textwidth}{0.6pt}\\[2pt]
\textbf{التعليمة:} يتضمن الاختبار خمس وضعيات إشكالية. اقرأ كل وضعية بعناية، ثم اختر التحليل أو التصرف الأنسب
من بين الاختيارات الخمسة؛ لكل وضعية إجابة صحيحة واحدة فقط.
اكتب \textbf{حرف} الإجابة الصحيحة في جدول الإجابات في آخر الاختبار.
لكل وضعية 4 نقاط (المجموع 20 نقطة).
\vspace{4pt}

\noindent
{\large\khatt\textbf{الوضعية 1}}\hfill{\small(الفصل 1)}\par\nopagebreak\vspace{3pt}
\setlength{\fboxsep}{7pt}\colorbox{gray!10}{\parbox{\dimexpr\linewidth-14pt}{تريد أستاذة نشاطًا تفاعليًا على الحاسوب: بطاقات يقلبها التلميذ ليرى إعراب كلمة. لا تعرف البرمجة، فاقترح زميل أن تطلب من أداة توليدية كتابة صفحة ويب بسيطة تنفّذ ذلك. وتخشى الأستاذة أن تكون هذه المهمة خارج قدرات الأدوات. ولديها أسبوع واحد لإعداد النشاط وتجربته قبل الحصة. ما التقدير الأدق؟}}\par\nopagebreak
\begin{itemize}[leftmargin=1.0cm, labelsep=6pt, itemsep=3pt, topsep=4pt, parsep=0pt]
  \item[\textbf{أ)}] الأدوات التوليدية تنتج النصوص والصور فقط، ولا تستطيع كتابة أي شيفرة برمجية، فالحل الوحيد أن تتعلم الأستاذة البرمجة أولًا، أو أن تكلّف مبرمجًا محترفًا بإنجاز المهمة، ولو استغرق ذلك أشهرًا.
  \item[\textbf{ب)}] الأدوات التوليدية تستطيع إنتاج شيفرة بسيطة من وصف لغوي، فتصف النشاط بدقة، وتجرّب الصفحة، وتطلب تعديلات عند الخلل، وتتحقق من صحة الإعراب في البطاقات قبل استعمالها مع التلاميذ.
  \item[\textbf{ج)}] تطلب الشيفرة وتنشرها مباشرة للتلاميذ دون تجربة، لأن الشيفرة التي تولّدها الأداة تعمل دائمًا من المحاولة الأولى، وإعراب البطاقات صحيح بالضرورة لأنه جزء من الشيفرة.
  \item[\textbf{د)}] الأدوات تكتب الشيفرة لكنها تشترط أن يكتب المستخدم طلبه بمصطلحات برمجية إنجليزية دقيقة، فلا يمكن للأستاذة وصف النشاط باللغة العربية، ولا الحصول على نتيجة قابلة للاستعمال، وهذا ما يعيق أغلب الأساتذة.
  \item[\textbf{هـ)}] تطلب الشيفرة وتجربها، لكنها لا تتحقق من الإعراب، لأن مسؤوليتها تقتصر على عمل الصفحة تقنيًا، أما المحتوى اللغوي فتتحمل الأداة مسؤوليته وحدها أمام التلاميذ والأولياء.
\end{itemize}
\par\vspace{6pt}
\noindent
{\large\khatt\textbf{الوضعية 2}}\hfill{\small(الفصل 2)}\par\nopagebreak\vspace{3pt}
\setlength{\fboxsep}{7pt}\colorbox{gray!10}{\parbox{\dimexpr\linewidth-14pt}{صحّحت أستاذة للأداة خطأً في إعراب «لا سيّما» داخل محادثة، فاعتذرت الأداة وصحّحت إجابتها. بعد أسبوع، في محادثة جديدة، تكرر الخطأ نفسه. استغربت الأستاذة لأنها «علّمت» الأداة من قبل، وسألت زملاءها عن السبب. وتريد أن تفهم السبب قبل أن تقرر هل تستمر في تصحيح الأداة. ما التفسير الأدق؟}}\par\nopagebreak
\begin{itemize}[leftmargin=1.0cm, labelsep=6pt, itemsep=3pt, topsep=4pt, parsep=0pt]
  \item[\textbf{أ)}] الاعتذار دليل على أن النموذج اقتنع وغيّر معارفه، وتكرار الخطأ سببه أن الأستاذة استعملت حسابًا مختلفًا، ولو استعملت الحساب نفسه لظهر التصحيح في كل محادثة جديدة.
  \item[\textbf{ب)}] النموذج تعلّم التصحيح وأضافه إلى معارفه، لكنه تراجع عنه لاحقًا لأن مستخدمين آخرين أعادوه إلى الخطأ، فالنموذج يعتمد على تصويت المستخدمين في كل معلومة يقدمها.
  \item[\textbf{ج)}] النموذج لا يتعلّم عادةً من محادثات المستخدم أثناء الاستعمال؛ التصحيح أثّر داخل تلك المحادثة فقط لأنه صار جزءًا من سياقها، أما معارف النموذج فتتغير بإعادة تدريبه لا بتصحيح فردي.
  \item[\textbf{د)}] النموذج لا يتعلّم من المحادثات، ولذلك كان اعتذاره وتصحيحه في المحادثة الأولى شكليين، ولم يتغير شيء حتى داخل تلك المحادثة، فالإجابات اللاحقة فيها بقيت تحمل الخطأ نفسه، وهذا ما يفسّر تكرار الخطأ بعد أسبوع.
  \item[\textbf{هـ)}] النموذج تعلّم التصحيح، لكنه نسيه بعد أسبوع لأن ذاكرته قصيرة كذاكرة الإنسان، ولو كررت الأستاذة التصحيح يوميًا لثبت نهائيًا في معارف النموذج عند كل المستخدمين، مثلما يثبت الحفظ عند التلاميذ بالتكرار.
\end{itemize}
\par\vspace{6pt}
\noindent
{\large\khatt\textbf{الوضعية 3}}\hfill{\small(الفصل 3)}\par\nopagebreak\vspace{3pt}
\setlength{\fboxsep}{7pt}\colorbox{gray!10}{\parbox{\dimexpr\linewidth-14pt}{طلب أستاذ صياغة رسالة إلى أولياء التلاميذ الضعاف في القراءة لدعوتهم إلى لقاء، فجاءت الرسالة رسمية جافة، فيها عبارات قد يفهمها الأولياء لومًا لهم، ومصطلحات تربوية غير مفهومة، ولم تحدد موعد اللقاء. وأولياء القسم متفاوتو المستوى التعليمي. وقد سبق أن اشتكى بعض الأولياء من رسائل مماثلة. أي طلب تعديل هو الأنسب؟}}\par\nopagebreak
\begin{itemize}[leftmargin=1.0cm, labelsep=6pt, itemsep=3pt, topsep=4pt, parsep=0pt]
  \item[\textbf{أ)}] «أعد صياغة الرسالة بأسلوب أدبي بليغ غني بالمحسنات البديعية»، لأن البلاغة تخفف وقع العبارات الجافة، وتجعل الرسالة أكثر تأثيرًا في كل الأولياء مهما كان مستواهم التعليمي.
  \item[\textbf{ب)}] «أعد صياغة الرسالة وأضف إليها نتائج كل تلميذ بالتفصيل»، لأن ذكر النتائج في الرسالة الجماعية يقنع الأولياء بالحضور، ويغني عن تعديل النبرة أو المصطلحات أو طول الرسالة.
  \item[\textbf{ج)}] «أعد صياغة الرسالة باللهجة العامية كاملة»، لأن العامية وحدها تضمن فهم كل الأولياء، وتلغي تلقائيًا أي نبرة رسمية أو لوم، دون حاجة إلى تحديد الطول أو المضمون أو موعد اللقاء.
  \item[\textbf{د)}] «أعد صياغة الرسالة بنبرة ودية دون مصطلحات تربوية»، ثم يرسلها مباشرة دون قراءتها، لأن تحديد النبرة كافٍ لضمان خلوّ الرسالة من أي عبارة قد تُفهم لومًا للأولياء.
  \item[\textbf{هـ)}] «أعد صياغة الرسالة بنبرة ودية مطمئنة لأولياء متفاوتي المستوى التعليمي، دون مصطلحات تربوية، مركّزًا على التعاون لا التقصير، في ثمانية أسطر، واختم بموعد اللقاء ومكانه.»
\end{itemize}
\par\vspace{6pt}
\noindent
{\large\khatt\textbf{الوضعية 4}}\hfill{\small(الفصل 5)}\par\nopagebreak\vspace{3pt}
\setlength{\fboxsep}{7pt}\colorbox{gray!10}{\parbox{\dimexpr\linewidth-14pt}{طلبت أستاذة من أداة استخراج الأفكار الأساسية والثانوية من نص حجاجي لتبني منه خريطة ذهنية لتلاميذها. جاءت القائمة سليمة الصياغة، لكن الأداة جعلت مثالًا توضيحيًا فكرة أساسية، ودمجت حجتين مستقلتين في فكرة واحدة. والخريطة ستُستعمل في المراجعة للاختبار. وستوزّع الخريطة مطبوعة على تلاميذ ثلاثة أقسام. ما التصرف الأنسب؟}}\par\nopagebreak
\begin{itemize}[leftmargin=1.0cm, labelsep=6pt, itemsep=3pt, topsep=4pt, parsep=0pt]
  \item[\textbf{أ)}] تبني الخريطة من القائمة كما هي، ثم تطلب من التلاميذ اكتشاف أخطائها دون تصحيحها مسبقًا، لأن الخريطة الذهنية أداة تقويم لا أداة تعلّم، فلا يضر أن تكون بنيتها خاطئة، فتكتشف الأستاذة بذلك مستوى فهمهم.
  \item[\textbf{ب)}] تعتمد القائمة كما هي، لأن الأداة تميّز الأفكار الأساسية من الثانوية أدق من القارئ البشري، وما تعدّه الأستاذة مثالًا قد يكون فكرة أساسية في نظر الكاتب لم تنتبه إليها، خاصة في النصوص الحجاجية.
  \item[\textbf{ج)}] تطلب من الأداة إعادة الاستخراج، وتعتمد النسخة الثانية دون مقارنة بالنص، لأن الأداة تتعلم من خطئها الأول داخل المحادثة، فتأتي المحاولة الثانية مطابقة لبنية النص بالضرورة.
  \item[\textbf{د)}] تراجع بنية النص بنفسها، فتعيد المثال تحت الحجة التي يخدمها، وتفصل الحجتين، ثم تبني الخريطة، لأن تمييز الأساسي من الثانوي حكم تحليلي يُتحقق منه في النص لا في صياغة الأداة.
  \item[\textbf{هـ)}] تتخلى عن الخريطة الذهنية وتكتفي بقائمة خطية، لأن الخطأ في ترتيب الأفكار يزول عندما تُعرض متتابعة دون تمييز بين الأساسي والثانوي، فلا حاجة إلى أي مراجعة للنص.
\end{itemize}
\par\vspace{6pt}
\noindent
{\large\khatt\textbf{الوضعية 5}}\hfill{\small(الفصل 6)}\par\nopagebreak\vspace{3pt}
\setlength{\fboxsep}{7pt}\colorbox{gray!10}{\parbox{\dimexpr\linewidth-14pt}{أنتج أستاذ بأداة توليد فيديو مقطعًا قصيرًا يمثّل سوق عكاظ، ليقدّم به درسًا عن الشعر الجاهلي. جاء المقطع جذابًا، لكن فيه أبنية حديثة، ولباسًا من عصور لاحقة، وشخصيات تحمل كتبًا مطبوعة. ويريد عرضه في الحصة الأولى من الوحدة. ويرى زميله أن هذه التفاصيل لا أهمية لها. ما الموقف الأنسب؟}}\par\nopagebreak
\begin{itemize}[leftmargin=1.0cm, labelsep=6pt, itemsep=3pt, topsep=4pt, parsep=0pt]
  \item[\textbf{أ)}] يعدّل الطلب بوصف تاريخي دقيق للمكان واللباس والأدوات، ويراجع المقطع مشهدًا مشهدًا، فإن بقيت مفارقات حذفها أو نبّه إليها، لأن المقطع الجذاب قد يرسّخ تصورًا تاريخيًا خاطئًا.
  \item[\textbf{ب)}] يعدّل الطلب بوصف تاريخي دقيق، ثم يعرض المقطع الجديد مباشرة دون مراجعة، لأن الوصف الدقيق يضمن وحده خلوّ المقطع من أي مفارقة تاريخية في المكان واللباس والأدوات.
  \item[\textbf{ج)}] يعرض المقطع كما هو، لأن هدف الدرس الشعر الجاهلي لا التاريخ، والتلاميذ لا ينتبهون إلى تفاصيل الأبنية واللباس، فلا تؤثر المفارقات في تعلّمهم أو في تصوّرهم للعصر الجاهلي.
  \item[\textbf{د)}] يعرض المقطع ويخبر التلاميذ أنه تمثيل دقيق لسوق عكاظ، لأن أدوات توليد الفيديو تعتمد على مصادر تاريخية موثقة، والمفارقات الظاهرة اجتهادات فنية مقبولة من الأداة.
  \item[\textbf{هـ)}] يتخلى عن الفيديو ويعرض صورة واحدة مولَّدة بالطلب نفسه، لأن الصور الثابتة لا تقع في المفارقات التاريخية أبدًا، بخلاف المقاطع المتحركة التي تكثر فيها هذه الأخطاء.
\end{itemize}
\par\vspace{6pt}
\begin{center}
{\khatt\textbf{جدول الإجابات}}\\[3pt]
\renewcommand{\arraystretch}{1.6}
\begin{tabular}{|>{\centering\arraybackslash}m{1.9cm}|*{5}{>{\centering\arraybackslash}m{1.2cm}|}}
\hline
\textbf{رقم السؤال} & \textbf{1} & \textbf{2} & \textbf{3} & \textbf{4} & \textbf{5} \\ \hline
\textbf{الإجابة} &  &  &  &  &  \\[8pt] \hline
\end{tabular}

\vspace{8pt}
\renewcommand{\arraystretch}{2.2}
\begin{tabular}{|>{\centering\arraybackslash}m{7.5cm}|>{\centering\arraybackslash}m{3.2cm}|>{\centering\arraybackslash}m{3.2cm}|}
\hline
\textbf{الاسم واللقب} & \textbf{الفوج} & \textbf{العلامة} \\ \hline
 & & \hfill /20 \\ \hline
\end{tabular}
\end{center}
\end{document}
